\chapter*{Introduction générale}
\markboth{Introduction générale}{}
\addcontentsline{toc}{chapter}{Introduction générale}
L'accès à un stage ou à une première expérience professionnelle met en relation deux besoins complémentaires. Les étudiants cherchent des missions compatibles avec leur formation, leurs projets et les compétences qu'ils souhaitent approfondir. Les entreprises souhaitent identifier des profils capables de contribuer à leurs activités, tout en tenant compte du caractère formateur du stage. La rencontre entre ces besoins dépend de la qualité des informations disponibles et de la manière dont elles sont comparées.

Les curriculum vitæ et les descriptions d'offres constituent des documents riches, mais peu homogènes. Une technologie peut être citée par son nom, par un framework associé ou par une expression propre à une langue. Les expériences sont décrites selon des formats différents et les compétences déclarées n'ont pas toutes le même niveau de preuve. Une plateforme qui ne fait que stocker ces documents laisse à l'utilisateur l'intégralité du travail de rapprochement. À l'inverse, une plateforme qui réduit ce travail à une note opaque peut donner une impression de précision sans fournir les éléments permettant d'apprécier le résultat.

Ce mémoire étudie la problématique suivante : \emph{comment concevoir une plateforme de mise en relation étudiants-entreprises qui automatise une partie de la comparaison entre profils et offres, tout en conservant des résultats explicables, une architecture évolutive et des limites identifiables ?} La réponse proposée prend la forme de SmartRecruit, développé dans le cadre du stage MDL/26/12 réalisé chez Best Solutions. Le projet associe la gestion des parcours de recrutement à un moteur de recommandation fondé sur des représentations textuelles et des connaissances explicites sur les compétences.

Le premier objectif est fonctionnel : regrouper les profils étudiants, les offres publiées, les candidatures et les décisions dans une application cohérente. Le deuxième objectif est technique : séparer les responsabilités de l'application web, de la persistance et du traitement de texte, de manière à pouvoir faire évoluer le moteur sans réécrire tous les parcours. Le troisième objectif est méthodologique : fournir une explication du score et une démarche d'évaluation qui permette de distinguer le comportement du logiciel de la pertinence réelle de ses recommandations.

Le choix initial s'oriente vers une méthode légère combinant TF-IDF, similarité cosinus et couverture de compétences. Cette méthode constitue une base de comparaison contrôlable : elle permet d'inspecter les termes, les synonymes et les pondérations utilisés. Elle ne résout cependant pas toutes les difficultés linguistiques. La négation, le niveau de maîtrise, les compétences seulement évoquées et les documents numérisés comme images restent des sources d'erreur. L'architecture prévoit donc une possibilité d'évolution, sans confondre cette possibilité avec une intégration déjà réalisée.

L'apport du travail ne se limite pas au calcul d'une note. Il concerne aussi l'articulation entre cette note et les processus applicatifs : à quel moment la calculer, quelles données lui transmettre, comment la conserver, et comment réagir lorsque le service chargé du calcul devient indisponible. L'évaluation montre notamment qu'un mécanisme de secours n'est utile que si son comportement est suffisamment cohérent avec le service principal. De même, la continuité d'affichage ne garantit pas la continuité des données lorsque plusieurs stockages indépendants sont utilisés.

Le mémoire est organisé en six chapitres. Le chapitre~\ref{chap:contexte} présente le contexte, la problématique et l'état de l'art. Le chapitre~\ref{chap:besoins} formalise les besoins et les critères d'acceptation. Le chapitre~\ref{chap:conception} expose l'architecture, le modèle de données et les interactions. Le chapitre~\ref{chap:realisation} décrit les éléments effectivement présents dans l'implémentation. Le chapitre~\ref{chap:moteur} détaille les traitements et les formules du moteur. Enfin, le chapitre~\ref{chap:validation} confronte les choix aux vérifications disponibles, discute les limites et propose une progression vers une version plus fiable. Les annexes réunissent les contrats d'échange, des scénarios synthétiques et les éléments de reproductibilité.
